\documentclass{article}
\usepackage{mathtools} 
\usepackage{fontspec}
\usepackage[UTF8]{ctex}
\usepackage{amsthm}
\usepackage{mdframed}
\usepackage{xcolor}
\usepackage{amssymb}
\usepackage{amsmath}


% 定义新的带灰色背景的说明环境 zremark
\newmdtheoremenv[
  backgroundcolor=gray!10,
  % 边框与背景一致，边框线会消失
  linecolor=gray!10
]{zremark}{说明}

% 通用矩阵命令: \flexmatrix{矩阵名}{元素符号}{行数}{列数}
\newcommand{\flexmatrix}[4]{
  \[
  #1 = \begin{pmatrix}
    #2_{11}     & #2_{12}     & \cdots & #2_{1#4}   \\
    #2_{21}     & #2_{22}     & \cdots & #2_{2#4}   \\
    \vdots      & \vdots      & \ddots & \vdots     \\
    #2_{#31}    & #2_{#32}    & \cdots & #2_{#3#4}
  \end{pmatrix}
  \]
}

% 简化版命令（默认矩阵名为A，元素符号为a）: \quickmatrix{行数}{列数}
\newcommand{\quickmatrix}[2]{\flexmatrix{A}{a}{#1}{#2}}

\begin{document}
\title{习题 12.3}
\author{张志聪}
\maketitle

\section*{1}

\begin{itemize}
  \item (3)
        \begin{align*}
          \sum \frac{(-1)^n}{n^{p + \frac{1}{n}}}
           & = \sum \frac{(-1)^n}{n^p \cdot n^{\frac{1}{n}}}             \\
           & = \sum (-1)^n \frac{1}{n^p} \cdot \frac{1}{n^{\frac{1}{n}}} \\
        \end{align*}

        \begin{itemize}
          \item (i) $p < 0$；
                \begin{align*}
                  \lim\limits_{n \to \infty} \frac{1}{n^p} \cdot \frac{1}{n^{\frac{1}{n}}}
                   & = +\infty
                \end{align*}
                故级数发散。

          \item (ii) $p = 0$；
                \begin{align*}
                  \lim\limits_{n \to \infty} \frac{1}{n^0} \cdot \frac{1}{n^{\frac{1}{n}}}
                   & = \lim\limits_{n \to \infty} 1 \cdot \frac{1}{n^{\frac{1}{n}}} \\
                   & = 1
                \end{align*}
                故级数发散。

          \item (iii) $1 < p$；
                \begin{align*}
                  \lim\limits_{n \to \infty} \frac{\frac{1}{n^p} \cdot n^{\frac{1}{n}}}{\frac{1}{n^p}}
                   & = \lim\limits_{n \to \infty} n^{\frac{1}{n}} \\
                   & = 1
                \end{align*}
                因为$\sum \frac{1}{n^p}$收敛，由比较原则，级数绝对收敛。

          \item (iv) $0 < p \leq 1$；
                \begin{align*}
                  \lim\limits_{n \to \infty} \frac{\frac{1}{n^p} \cdot n^{\frac{1}{n}}}{\frac{1}{n^p}}
                   & = \lim\limits_{n \to \infty} n^{\frac{1}{n}} \\
                   & = 1
                \end{align*}
                因为$\sum \frac{1}{n^p}$发散，由比较原则，级数不绝对收敛。

                考虑条件收敛，对于交错级数$\sum (-1) \frac{1}{n^p}$，由莱布尼茨判别法可知其收敛。
                又有
                \begin{align*}
                  (\sqrt[x]{x})'
                   & = (e^{\ln\, \sqrt[x]{x}})'                                     \\
                   & = (e^{\frac{1}{x} \ln\, x})'                                   \\
                   & = e^{\frac{1}{x} \ln\, x} \left[\frac{1 - \ln\, x}{x^2}\right] \\
                   & = \sqrt[x]{x} \left[\frac{1 - \ln\, x}{x^2}\right]
                \end{align*}
                所以，$x > e$时，函数$\sqrt[x]{x}$单减，由归结原理可知，$n > 3$时，$\sqrt[n]{n}$单调且有界，
                因为改变级数的有限项并不影响原有级数的敛散性，因此利用阿贝尔判别法，
                级数收敛。
        \end{itemize}

\end{itemize}

\end{document}